\documentclass[a4paper]{report}
% Shared preamble for the English probability and statistics documents.
% Based on the original course preamble. Compile with pdfLaTeX.
\usepackage[top=1.5cm,bottom=2cm,left=2cm,right=2cm]{geometry}
\usepackage[T1]{fontenc}
\usepackage{lmodern}
\usepackage[utf8]{inputenc}
\usepackage[english]{babel}

\usepackage{amsmath}
\usepackage{amssymb}
\usepackage{amsthm}
\usepackage{mathrsfs}
\usepackage{mathtools}
\usepackage{etoolbox}
\usepackage{tikz}
\usetikzlibrary{shapes.multipart,arrows.meta}
\usepackage{xcolor}
\usepackage{enumerate}
\usepackage[shortlabels]{enumitem}
\usepackage{hyperref}
\usepackage{version}
\usepackage{bbold}
\usepackage{fontawesome5}
\usepackage{fancyhdr}
\lfoot{}
\usepackage{natbib}
\usepackage{graphicx}
\usepackage{pdfpages}
\usepackage{authblk}
    
\usepackage{blindtext}
\usepackage{multicol}
\usepackage{collcell}
\usepackage{wrapfig}
\usepackage{blkarray}
\usepackage{cancel}
\usepackage{appendix}
\usepackage{float}
\makeatletter
\input{tex/assets/slashbox.sty}
\makeatother
\hypersetup{unicode,colorlinks,citecolor=red,filecolor=black,
  linkcolor=red,urlcolor=cyan}
\setlength{\columnsep}{1cm}
\pagenumbering{gobble}

\def\numtext#1{%
  \ifcase#1\or one\or two\or three\or four\or five\or six\or seven\or eight\or nine\or ten\or
  eleven\or twelve\or thirteen\or fourteen\or fifteen\or sixteen\or seventeen\or eighteen\or nineteen\or twenty\or Lots\fi}

\newcommand{\Byte}[1]{%
  \def\nodecontents{}%
  \foreach[count=\l] \k in {#1} {\xappto{\nodecontents}{\noexpand\nodepart{\numtext{\l}}\k}}%
  \node[rectangle split,rectangle split parts=8,rectangle split horizontal,draw] at (2,2){\nodecontents};}

\newcommand{\D}{{\mathrm{d}}}
\newcommand{\cE}{{\mathcal{E}}}
\newcommand{\cG}{{\mathcal{G}}}
\newcommand{\cX}{{\mathcal{X}}}
\newcommand{\cY}{{\mathcal{Y}}}
\newcommand{\cB}{{\mathcal{B}}}
\newcommand{\cF}{{\mathcal{F}}}
\newcommand{\cO}{{\mathcal{O}}}
\newcommand{\cS}{{\mathcal{S}}}
\newcommand{\cT}{{\mathcal{T}}}
\newcommand{\given}{\, \big| \,}
\newcommand{\cL}{{\mathscr{L}}}
\newcommand{\cD}{{\mathcal{D}}}
\newcommand{\cA}{{\mathscr{A}}}
\newcommand{\cP}{{\mathscr{P}}}
\newcommand{\cR}{{\mathscr{R}}}
\newcommand{\A}{{\mathcal A}}
\newcommand{\one}{\mathbb{1}}
\newcommand{\mathbbm}[1]{\mathbb{#1}}

\providecommand{\C}{}
\renewcommand{\C}{\mathbb C}
\newcommand{\R}{\mathbb R}
\newcommand{\Q}{\mathbb Q}
\newcommand{\N}{\mathbb N}
\newcommand{\T}{\mathbb T}
\newcommand{\F}{\mathbb F}
\renewcommand{\P}{\mathbb P}
\newcommand{\Z}{\mathbb Z}
\newcommand{\E}{\mathbb {E}}
\newcommand{\Pbb}[1]{\P\left( #1\right)}
\newcommand{\Ebb}[1]{\E\left[ #1\right]}

\newcommand{\lb}{\left (}
\newcommand{\rb}{\right )}
\newcommand{\lbb}{\left [}
\newcommand{\rbb}{\right ]}
\newcommand{\labs}{\left |}
\newcommand{\rabs}{\right |}
\newcommand{\lbrb}[1]{\lb #1 \rb}
\newcommand{\lbbrbb}[1]{\lbb#1\rbb}
\newcommand{\labsrabs}[1]{\labs#1\rabs}
\newcommand{\lbbrb}[1]{\lbb#1\rb}
\newcommand{\lbrbb}[1]{\lb#1\rbb}
\newcommand{\langlerangle}[1]{\langle#1\rangle}
\newcommand{\lbcurly}{\{}
\newcommand{\rbcurly}{\}}
\newcommand{\curly}[1]{\lbcurly#1\rbcurly}
\newcommand{\abs}{\labsrabs}
\newcommand{\ind}[1]{1_{#1}}
\newcommand{\indep}{\perp \!\!\! \perp}
\renewcommand{\d}{\, \mathrm{d}}

\newtheorem{theorem}{Theorem}
\newtheorem{problemCnt}{Lemma}
\theoremstyle{definition}
\newtheorem{problem}[problemCnt]{Problem}
\newtheorem{problembook}[problemCnt]{\faBook \,  Problem}
\newtheorem{problempencil}[problemCnt]{\faPencil* \,  Problem}
\newtheorem{problemstar}[problemCnt]{\faStar\,  Problem}
\newtheorem{problemstarpencil}[problemCnt]{\faStar \, \faPencil* \,   Problem}
\newtheorem{probleminteresting}[problemCnt]{\faSmile[regular] \,   Problem}
\newtheorem{solution}{Solution}

\newcommand{\fwcell}[1]{\makebox[\arraystretch\normalbaselineskip]{$#1$}}
\newenvironment{enumeratei}{\begin{enumerate}[(i)]}{\end{enumerate}}
\newenvironment{enumeratea}{\begin{enumerate}[(a)]}{\end{enumerate}}
\newenvironment{enumeraten}{\begin{enumerate}[\upshape 1.]}{\end{enumerate}}

% Compatibility helpers needed by a few files in the combined archive.
\AtBeginDocument{\noindent{\footnotesize English translation prepared with ChatGPT (OpenAI, GPT-6), 26 October 2026, for translation of the Bulgarian text into English and LaTeX formatting.
The Bulgarian originals remain the reference; translation may contain errors.}\par\smallskip}
\newcommand{\MissingGraphic}[1]{%
  \begin{center}\fbox{\parbox{0.82\linewidth}{\centering Missing source image: \texttt{#1}}}\end{center}}



\begin{document}
\noindent Stochastic Empirical Methods, SE Year 3, 2023/2024 \hfill Bonus homework
\hrulefill
\begin{it}This assignment awards no points, but illustrates some types of argument required in the theory examination. Enjoy the work! Work on $(\Omega,\mathcal A,\P)$. Write $\overline A=\Omega\setminus A$ and $\mathbb N_0=\mathbb N\cup\{0\}$.\end{it}
\begin{problem}
For $A,B\subset\Omega$, define $A\mathbin\Delta B=(A\cap\overline B)\cup(\overline A\cap B)$. Express $\mathbbm1_{(A\Delta B)\Delta C}$ using $\mathbbm1_A,\mathbbm1_B,\mathbbm1_C$.
\end{problem}
\begin{problem}
Prove for $A,B\subset\Omega$ with $\P(B)\ne0$ that $\P(A\mid B)\geq[\P(A)+\P(B)-1]/\P(B)$.
\end{problem}
\begin{problem}
Pool samples from $n$ people for a COVID test. If negative, declare everyone healthy; if positive, retest each individually. Four equally common blood groups each represent a quarter of the population, with infection proportions 10\%, 1\%, 5\% and 4\%, respectively. Each person's blood group is independent of the others, with these population proportions. The pooled test costs 1 BGN and each individual test 1.25 BGN.
\begin{itemize}
\item Model the expected cost $p_c(n)$ per tested person.
\item Treat $n$ as continuous $x$. Derive, without solving, an equation for $x^*$ satisfying \[p_c(x^*)=\min_{x>0}p_c(x).\]
\item Optionally use a computer to find the optimal integer $n^*$ and minimum cost.
\end{itemize}
\end{problem}
\begin{problem}
Let $X$ be $\mathbb N_0$-valued and $H_X(\lambda)=\E[e^{-\lambda X}]$ for $\lambda\geq0$.
\begin{itemize}
\item Prove $H_X(\lambda)\leq1$ and, by formally differentiating under the expectation, express $\E X$ and $\E X^2$.
\item Express the generating function $s_X$ using $H_X$.
\item For $X\sim\mathrm{Bin}(n,p)$, find $H_X$ and use it to obtain $\E X$, $\E X^2$.
\end{itemize}
\end{problem}
\begin{problem}
A power plant's daily operation depends on $Y=|X-1|$, the deviation of a parameter with $\E X=1$, $DX=0.01$. Because large deviations pose an environmental hazard, a regulator imposes a fine $g(n)$ BGN if $Y\in(n,n+1]$, $n\geq2$.
\begin{itemize}
\item Find an upper bound for $f(a)=\P(Y>a)$.
\item Independent experts claim that for $g(n)=n^{3/2}$, $\lim_{n\to\infty}g(n)\P(Y\in[n,n+1))=0$, making the fines insufficiently effective. Is their claim about the limit true?
\end{itemize}
\end{problem}
\begin{problem}
Let $X_1,X_2,\ldots$ be independent Poisson variables with $\P(X_1=k)=e^{-\lambda}\lambda^k/k!$, $k\geq0$. Model a birth process by $Y_n=(1+X_n)Y_{n-1}$, $Y_0=1$, and a birth-and-death process by $Z_n=X_nZ_{n-1}$, $Z_0=1$.
\begin{itemize}
\item Define almost sure convergence of random variables $\xi_n$ to $\xi$.
\item Does $n^{-1}\log Y_n$ have a limit, and what is it? What type of convergence holds? Does $\E Y_n$ grow exponentially?
\item Calculate $\P(Z_n\ne0)$ and prove that $Z_n\to0$ almost surely, so extinction occurs almost surely. Does $\E Z_n$ grow exponentially when $\lambda>1$?
\end{itemize}
\end{problem}
\begin{problem}
Roll a six-sided die $3\times10^{12}$ times and let $X$ count ones or threes. A struggling student is asked to approximate:
\begin{itemize}
\item $\P(X>10^{12})$;
\item $\P(X>10^{12}+10^3)$;
\item $\P(X>10^{12}+10^9)$.
\end{itemize}
The student answers ``fifty-fifty'', or $1/2$, to all three. How many answers, and which, are correct? Justify mathematically.
\end{problem}
\begin{problem}
The duration of acquired immunity is exponential $X$ with rate $\lambda$. Immunity strength is a continuous variable $Y$ with density $f_Y(y)=\beta y^{\beta-1}$, $0<y\leq1$, $\beta>0$. Given $n$ observations $(x_1,\ldots,x_n)$ and $(y_1,\ldots,y_n)$, find maximum likelihood estimators of $\lambda,\beta$. Are they consistent?
\end{problem}
\begin{problem}
In a phase 3 vaccine trial, $n=10000$ people are tested and 9000 acquire immunity. Model the immune count as $X\sim\mathrm{Bin}(10000,p)$. Construct a symmetric critical region for $H_0:p=0.9$ with type I error $\alpha=0.01$. Is there reason to regard this region as close to optimal?
\end{problem}
\end{document}

